% References for paper.tex (\input before \appendix). Author-year, alphabetical. Details come from
% studies/literature.md, which records what was read and how; where it gives no title or venue, the entry
% gives only what is known (no guessed details); a description in brackets stands in for an unchecked title. A new citation in the paper gets a line here.
\section*{References}\addcontentsline{toc}{section}{References}
\label{sec:references}
\begingroup\small\setlength{\parindent}{0pt}\setlength{\parskip}{2pt}
\newcommand{\bib}[1]{\par\hangindent=1.5em\hangafter=1 #1}
\bib{Aitken, A. C. (1935). On least squares and linear combination of observations. \emph{Proc.\ R.\ Soc.\ Edinb.} 55, 42--48.}
\bib{Andblom, M. (2023). Generalized Bühlmann--Straub credibility theory for correlated data. Master thesis 2023:2, Stockholm University.}
\bib{Beck, N. and Katz, J. N. (1995). What to do (and not to do) with time-series cross-section data. \emph{Am.\ Polit.\ Sci.\ Rev.} 89, 634--647.}
\bib{Bell, W. R., Datta, G. S. and Ghosh, M. (2013). Benchmarking small area estimators. \emph{Biometrika} 100.}
\bib{Bloomfield, P. and Watson, G. S. (1975). The inefficiency of least squares. \emph{Biometrika} 62.}
\bib{Bodnar, T., Parolya, N. and Schmid, W. (2018). Estimation of the global minimum variance portfolio in high dimensions. \emph{Eur.\ J. Oper.\ Res.} 266.}
\bib{Bühlmann, H. (1967). Experience rating and credibility. \emph{ASTIN Bull.} 4(3).}
\bib{Bühlmann, H. (1974). A comparison of three credibility formulae using multidimensional techniques. \emph{ASTIN Bull.} 7(3), 203--207.}
\bib{Bühlmann, H. and Gisler, A. (2005). \emph{A Course in Credibility Theory and its Applications.} Springer.}
\bib{Bühlmann, H. and Straub, E. (1970). Glaubwürdigkeit für Schadensätze. \emph{Mitt.\ Ver.\ Schweiz.\ Vers.-Math.} 70.}
\bib{Cossette and Luong (2003). [Weighted least squares for regression credibility with errors correlated within a risk; title not checked.]}
\bib{Dannenburg, D. R., Kaas, R. and Goovaerts, M. J. (1996). \emph{Practical Actuarial Credibility Models.} IAE, Amsterdam.}
\bib{Das, K., Jiang, J. and Rao, J. N. K. (2004). Mean squared error of empirical predictor. \emph{Ann.\ Stat.} 32.}
\bib{Datta et al.\ (2011). [Bayesian benchmarking in small-area estimation; title not checked.]}
\bib{Ebrahimzadeh, Ibrahim, Jemain and Kilicman (2012). Unbiased estimation of structural parameters in credibility models with dependence induced by common effects. \emph{Bull.\ Malays.\ Math.\ Sci.\ Soc.} 35(4).}
\bib{Ebrahimzadeh et al.\ (2013). Claim dependence with common effects in credibility models with error uniform dependence. \emph{World Appl.\ Sci.\ J.} 27(12), 1551--1560.}
\bib{Fisher, W. D. (1958). On grouping for maximum homogeneity. \emph{J. Am.\ Stat.\ Assoc.} 53, 789--798.}
\bib{Frahm, G. and Memmel, C. (2010). Dominating estimators for minimum-variance portfolios. \emph{J. Econometrics} 159.}
\bib{Frees, E. W., Young, V. R. and Luo, Y. (1999). A longitudinal data analysis interpretation of credibility models. \emph{Insur.\ Math.\ Econ.} 24, 229--247.}
\bib{Fung and Xu (2009). [ML and REML for the structural parameters of regression credibility; title not checked.] \emph{Variance}.}
\bib{Furrer, R., Genton, M. G. and Nychka, D. (2006). Covariance tapering for interpolation of large spatial datasets. \emph{J. Comput.\ Graph.\ Stat.} 15(3), 502--523.}
\bib{Gerber, H. U. and Jones, D. A. (1975). Credibility formulas of the updating type. \emph{Trans.\ Soc.\ Actuar.} 27.}
\bib{Gisler, A. and Wüthrich, M. V. (2008). Credibility for the chain ladder reserving method. \emph{ASTIN Bull.} 38(2).}
\bib{Goldberger, A. S. (1962). Best linear unbiased prediction in the generalized linear regression model. \emph{J. Am.\ Stat.\ Assoc.} 57.}
\bib{Hachemeister, C. A. (1975). Credibility for regression models with application to trend. In \emph{Credibility: Theory and Applications}, Academic Press.}
\bib{Happ, Maier and Merz (2011). [Multivariate Bühlmann--Straub for correlated run-off triangles; title not checked.]}
\bib{Harville, D. A. (1976). Extension of the Gauss--Markov theorem to include the estimation of random effects. \emph{Ann.\ Stat.} 4(2), 384--395.}
\bib{Henderson, C. R. (1963). Selection index and expected genetic advance. In \emph{Statistical Genetics and Plant Breeding}, NAS--NRC 982.}
\bib{Henderson, C. R. (1975). Best linear unbiased estimation and prediction under a selection model. \emph{Biometrics} 31, 423--447.}
\bib{Huang and Wu (2015). Credibility models with dependence structure over risks and time horizon. \emph{J. Ind.\ Manag.\ Optim.} 11(2), 365--380.}
\bib{Jewell, W. S. (1973). [Multidimensional credibility; unpublished report, title not checked.]}
\bib{Jewell, W. S. (1974). Credible means are exact Bayesian for exponential families. \emph{ASTIN Bull.} 8(1), 77--90.}
\bib{Jewell, W. S. (1975). [The hierarchical credibility model; title not checked.]}
\bib{Jiang, J., Lahiri, P. and Wan, S.-M. (2002). A unified jackknife theory for empirical best prediction with M-estimation. \emph{Ann.\ Stat.} 30.}
\bib{Jiang, Nguyen and Rao (2011). Best predictive small area estimation. \emph{J. Am.\ Stat.\ Assoc.} 106(494), 732--745.}
\bib{Kackar, R. N. and Harville, D. A. (1981). Unbiasedness of two-stage estimation and prediction procedures for mixed linear models. \emph{Comm.\ Stat.\ A} 10.}
\bib{Kackar, R. N. and Harville, D. A. (1984). Approximations for standard errors of estimators of fixed and random effects in mixed linear models. \emph{J. Am.\ Stat.\ Assoc.} 79.}
\bib{Klugman, S. A. (1992). \emph{Bayesian Statistics in Actuarial Science.} Kluwer.}
\bib{Ledoit, O. and Wolf, M. (2004). Honey, I shrunk the sample covariance matrix. \emph{J. Portf.\ Manag.} 30(4).}
\bib{Ledoit, O. and Wolf, M. (2017). Nonlinear shrinkage of the covariance matrix for portfolio selection: Markowitz meets Goldilocks. \emph{Rev.\ Financ.\ Stud.} 30(12).}
\bib{Lo, Fung and Zhu (2006). [GEE estimators of the structural parameters of regression credibility; title not checked.] \emph{ASTIN Bull.}}
\bib{Mack, T. (1993). Distribution-free calculation of the standard error of chain ladder reserve estimates. \emph{ASTIN Bull.} 23(2).}
\bib{Morris, J. (2026). Logistic credibility with temporal decay. arXiv:2606.08692.}
\bib{Parodi, P. and Bonche, S. (2010). Uncertainty-based credibility and its applications. \emph{Variance}.}
\bib{Pfeffermann, D. and Barnard, C. H. (1991). Some new estimators for small-area means with application to the assessment of farmland values. \emph{J. Bus.\ Econ.\ Stat.} 9.}
\bib{Poon (2013). Spatial cross-sectional credibility models with general dependence structure among risks. MSc project, Simon Fraser University.}
\bib{Poon and Lu (2015). A spatial cross-sectional credibility model with dependence among risks. \emph{N. Am.\ Actuar.\ J.} 19(4).}
\bib{Prasad, N. G. N. and Rao, J. N. K. (1990). The estimation of the mean squared error of small-area estimators. \emph{J. Am.\ Stat.\ Assoc.} 85.}
\bib{Reluga, K., Salvati, N., Ranjbar, S., Tzavidis, N. and van der Laan, M. (2026). Small area learning. arXiv:2610.15953.}
\bib{Schnapp et al.\ (2000). [Crop-insurance (MPCI) ratemaking with spatially correlated effects; title not checked.] \emph{CAS Forum}.}
\bib{Stein, C. M. (1981). Estimation of the mean of a multivariate normal distribution. \emph{Ann.\ Stat.} 9.}
\bib{Stein, M. L. (1999). \emph{Interpolation of Spatial Data: Some Theory for Kriging.} Springer.}
\bib{Stukel and Rao (1999). On small-area estimation under two-fold nested error regression models. \emph{J. Stat.\ Plan.\ Inference} 78, 131--147.}
\bib{Styan (1983). [Survey of the Kantorovich and Wielandt inequalities for the efficiency of least squares; title not checked.]}
\bib{Taylor, G. C. (1979). Credibility analysis of a general hierarchical model. \emph{Scand.\ Actuar.\ J.} 1979(1), 1--12.}
\bib{van Rijn, Weber and Merz (2026). Credibility clustering for experience rating: an automated approach to risk group allocation in group life insurance. arXiv:2611.07110.}
\bib{Wang, J., Fuller, W. A. and Qu, Y. (2008). Small area estimation under a restriction. \emph{Survey Methodology} 34.}
\bib{Wen, L. and Wu, X. (2011). The credibility estimator with general dependence structure over risks. \emph{Comm.\ Stat.\ Theory Methods} 40(10), 1893--1910.}
\bib{Wen, L., Wu, X. and Zhou, X. (2009). The credibility premiums for models with dependence induced by common effects. \emph{Insur.\ Math.\ Econ.} 44(1).}
\bib{Yang and Wang (2024). [REML and linear-Bayes estimators for the Bühlmann model with correlated risks; title not checked.] \emph{Comm.\ Stat.} 53(8).}
\bib{Yeo, K. L. and Valdez, E. A. (2006). Claim dependence with common effects in credibility models. \emph{Insur.\ Math.\ Econ.} 38(3).}
\bib{Zellner, A. (1962). An efficient method of estimating seemingly unrelated regressions and tests for aggregation bias. \emph{J. Am.\ Stat.\ Assoc.} 57.}
\bib{Zhang and Chen (2018). Credibility estimators with dependence structure over risks and time under balanced loss function. \emph{Stat.\ Neerl.} 72(2), 157--173.}
\bib{Zhang, Cui and Chen (2017). Multidimensional credibility estimators with random common effects and time effects. \emph{J. Syst.\ Sci.\ Complex.} 30, 1107--1120.}
\par\endgroup
